\documentclass[11pt,a4paper]{article}
\usepackage{fontspec,amsmath,amssymb,amsthm,microtype}
\setmainfont{lmroman10-regular.otf}[BoldFont=lmroman10-bold.otf,ItalicFont=lmroman10-italic.otf,BoldItalicFont=lmroman10-bolditalic.otf]
\setmonofont{lmmono10-regular.otf}[Scale=MatchLowercase]
\usepackage[margin=25mm]{geometry}
\usepackage[unicode,hidelinks]{hyperref}
\setlength{\emergencystretch}{3em}
\setlength{\parskip}{0.4em}
\setlength{\parindent}{0pt}
\allowdisplaybreaks[2]
\raggedbottom
\title{\texorpdfstring{The dimension-\allowbreak{}sharp equilibrium parameter concentration}{The dimension-sharp equilibrium parameter concentration}}
\author{LangevinLean mathematical exposition}
\date{}
\begin{document}\maketitle
\noindent Source module: \texttt{IPLA/\allowbreak{}Concentration.\allowbreak{}lean}. The accompanying HTML pairs this mathematical exposition with the exact Lean statements and records the scope of the formalization.


The exact parameter marginal has density proportional to \(M(\theta)^N=e^{-N\Phi(\theta)}\). Its potential is \(Nm\)-strongly convex in dimension \(d_\theta\), giving the equilibrium radius \(\sqrt{2d_\theta/(Nm)}\) without a separate concentration hypothesis.


\section{\texorpdfstring{Identify the scaled marginal potential and its Gibbs measure}{Identify the scaled marginal potential and its Gibbs measure}}


Multiply the strong-convexity inequality for \(\Phi\) by the nonnegative scalar \(N\). Every affine term and the quadratic penalty scale by \(N\), so \(N\Phi\) is \(Nm\)-strongly convex. Since \(M>0\),

\[e^{-N\Phi(\theta)}=e^{N\log M(\theta)}=M(\theta)^N.\]

The finite-real representation of the original extended marginal identifies this weight exactly with the one defining \texttt{marginal\allowbreak{}Target}. Its normalized law is therefore the shared Gibbs measure of \(N\Phi\).




\textbf{Correspondence with the displayed declarations.}


\begin{itemize}

\item \texttt{scaled\_\allowbreak{}marginal\allowbreak{}Potential\_\allowbreak{}strong\allowbreak{}Convex}: Multiplying the actual marginal potential by \texttt{N} multiplies its proved strong-convexity constant by \texttt{N}.
\item \texttt{marginal\allowbreak{}Target\_\allowbreak{}eq\_\allowbreak{}gibbs\allowbreak{}Measure}: The normalized powered marginal density is exactly the Gibbs law of the proved \texttt{N}-scaled negative log marginal.
\end{itemize}


{\raggedright\noindent Lean declarations: \texttt{scaled\_\allowbreak{}marginal\allowbreak{}Potential\_\allowbreak{}strong\allowbreak{}Convex}, \texttt{marginal\allowbreak{}Target\_\allowbreak{}eq\_\allowbreak{}gibbs\allowbreak{}Measure}.\par}

\section{\texorpdfstring{A dilation estimate gives the sharp parameter moment}{A dilation estimate gives the sharp parameter moment}}


A marginal maximizer \(\theta_*\) minimizes \(\Phi\), hence also \(N\Phi\). The shared continuous strongly convex Gibbs theorem, which needs continuity and a mode rather than differentiability of \(\Phi\), bounds the squared distance by twice dimension over convexity:

\[\int\|\theta-\theta_*\|^2d\pi_N^\theta(\theta)\le\frac{2d_\theta}{Nm}.\]

Apply it to the continuous potential \(N\Phi\), use the proved \(Nm\)-strong convexity, and rewrite its Gibbs law by the preceding identity. Only the parameter dimension enters; latent dimension does not enlarge this equilibrium radius.




\textbf{Correspondence with the displayed declarations.}


\begin{itemize}

\item \texttt{marginal\allowbreak{}Target\_\allowbreak{}second\allowbreak{}Moment\_\allowbreak{}le}: Exact finite second-moment bound about any parameter maximizing the actual latent marginal. Its existence and uniqueness were proved separately.
\end{itemize}


{\raggedright\noindent Lean declarations: \texttt{marginal\allowbreak{}Target\_\allowbreak{}second\allowbreak{}Moment\_\allowbreak{}le}.\par}

\section{\texorpdfstring{Transfer the marginal moment to the lifted Gibbs target}{Transfer the marginal moment to the lifted Gibbs target}}


The exact projection identity \(p_{\theta\#}\Pi_N=\pi_N^\theta\) and the measurable pushforward integral formula yield

\[\int\|p_\theta z-\theta_*\|^2d\Pi_N(z)\le\frac{2d_\theta}{Nm}.\]

The projected Gibbs variable is already in \(L^2\), and subtracting the constant \(\theta_*\) preserves this. The shared projected-RMS theorem converts the real moment inequality to the extended RMS bound

\[\operatorname{RMS}_{\Pi_N}(p_\theta-\theta_*)\le\sqrt{2d_\theta/(Nm)}.\]

The radius is nonnegative and its square equals the moment bound, so the square-root step is legitimate. This term needs joint differentiability and strong convexity, with no gradient-Lipschitz premise.




Powering the marginal density by \(N\) multiplies its negative log potential by \(N\), so its modulus becomes \(Nm\). The actual marginal maximizer minimizes that potential. The dilation concentration theorem used in this module bounds its parameter second moment by \(2d_\theta/(Nm)\), and the exact parameter pushforward identity gives the same moment under the lifted Gibbs law. The final square-root estimate uses the proved finite second moment to interpret projected \(L^2\) error as an ordinary integral. The stronger \(d_\theta/(Nm)\) constant is separately proved in the proximal sharp-concentration extension.




\textbf{Correspondence with the displayed declarations.}


\begin{itemize}

\item \texttt{gibbs\allowbreak{}Target\_\allowbreak{}parameter\_\allowbreak{}second\allowbreak{}Moment\_\allowbreak{}le}: Transfer the proved marginal moment bound to the actual scaled IPLA Gibbs measure using the exact parameter-marginal identity.
\item \texttt{gibbs\allowbreak{}Target\_\allowbreak{}parameter\_\allowbreak{}concentration}: The equilibrium concentration input in IPLA Theorem 1 is a proved theorem of the original joint assumptions. No gradient-Lipschitz or marginal assumption is needed for this statistical term.
\end{itemize}


{\raggedright\noindent Lean declarations: \texttt{gibbs\allowbreak{}Target\_\allowbreak{}parameter\_\allowbreak{}second\allowbreak{}Moment\_\allowbreak{}le}, \texttt{gibbs\allowbreak{}Target\_\allowbreak{}parameter\_\allowbreak{}concentration}.\par}

\end{document}
